% GENERATED FILE. Edit canonical lessons, then run npm run generate:books.
% canonical-source-hash: fnv1a64:e80396ea043013d5
% canonical-lessons: ZH-C450-haixian, ZH-C450-tedian, ZH-W450-te-special, ZH-C450-teyi, ZH-C450-zhengzong, ZH-C450-zhaopaicai

\chapter{Local Food}
\label{ch:zh-local-food}
\hlchaptermodality{450}

\begin{chapteropening}
\textbf{By the end of this chapter:} \emph{I can talk about seafood, what makes the food special, going somewhere on purpose, authentic food and a signature dish.}

Complete the last lesson of chapter 450: I can talk about seafood, what makes the food special, going somewhere on purpose, authentic food and a signature dish.
\end{chapteropening}

\section[hǎixiān]{hǎixiān — seafood}

\label{lesson:ZH-C450-haixian}



\begin{quote}
Before the new one: say the Mandarin for light and mild (of food), then the Mandarin for greasy (of food).
\end{quote}

\subsection*{You'll want to know: hǎixiān}
\emph{hǎixiān} — \textquotedblleft{}seafood\textquotedblright{}.

\begin{quote}
\emph{zhè jiā fànguǎn de hǎixiān hěn xīnxiān} — The seafood at this restaurant is very fresh.
\end{quote}

Its characters come later; for now it is a word for the ear.

\begin{cousinweb}
\emph{hǎixiān} is built from \zh{海} 'the sea' and \emph{xiān} 'fresh', both inherited from Middle Chinese, \emph{xojX} and \emph{sjen} (Guangyun, in Baxter's transcription), which glosses the second 'fresh' and 'clean'.
\end{cousinweb}

\begin{grammarlens}[title={What you've built}]
One more word for getting around.
\end{grammarlens}

\subsection*{Your turn}
\begin{itemize}
\raggedright
  \item \emph{Say it:} \emph{hǎixiān}
  \item \emph{Say it:} \emph{hǎixiān}, once more
  \item \emph{Say it:} \emph{zhè jiā fànguǎn de hǎixiān hěn xīnxiān} — The seafood at this restaurant is very fresh.
  \item \emph{Recall:} say the Mandarin for light and mild (of food), then the Mandarin for greasy (of food), then say \emph{hǎixiān} again
\end{itemize}

\begin{tcolorbox}[breakable,colback=teal!4,colframe=teal!35!black,title={Before you move on}]
What does hǎixiān mean? (\textquotedblleft{}Seafood\textquotedblright{}.) Say it once more.
\end{tcolorbox}
\section[\texorpdfstring{\zh{特点} (tèdiǎn)}{特点 (tèdiǎn)}]{\zh{特点} (tèdiǎn) — characteristic; a distinguishing trait}

\label{lesson:ZH-C450-tedian}



\begin{quote}
Before the new one: say the Mandarin for greasy (of food), then the Mandarin for seafood.
\end{quote}

\subsection*{You'll want to know: \zh{特点}}
\textbf{\zh{特点}} — \emph{tèdiǎn} — \textquotedblleft{}characteristic; a distinguishing trait\textquotedblright{}.

\begin{quote}
\textbf{\zh{这儿的菜有什么特点}} \emph{zhèr de cài yǒu shénme tèdiǎn} — What's special about the food here?
\end{quote}

The new character, in pieces: \zh{特} — ten strokes in two pieces: on the left the cow piece, four strokes; on the right the earth piece of \zh{地} over the right piece of \zh{对}. The next lesson puts it in your hand.

\begin{cousinweb}
\textbf{\zh{特点}} (\emph{tèdiǎn}) is built from \zh{特} 'single, special' and \zh{点} 'a point', both inherited from Middle Chinese, \emph{dok} and \emph{temX} (Guangyun, in Baxter's transcription), which glosses the first as a bull and 'alone' and the second as a dot made with the brush.
\end{cousinweb}

\begin{grammarlens}[title={What you've built}]
One more word for getting around.
\end{grammarlens}

\subsection*{Your turn}
\begin{itemize}
\raggedright
  \item \emph{Say it:} \emph{tèdiǎn}
  \item \emph{Say it:} \emph{tèdiǎn}, once more
  \item \emph{Say it:} \emph{zhèr de cài yǒu shénme tèdiǎn} — What's special about the food here?
  \item \emph{Recall:} say the Mandarin for greasy (of food), then the Mandarin for seafood, then say \emph{tèdiǎn} again
\end{itemize}

\begin{tcolorbox}[breakable,colback=teal!4,colframe=teal!35!black,title={Before you move on}]
What does \zh{特点} mean? (\textquotedblleft{}Characteristic; a distinguishing trait\textquotedblright{}.) Say it once more.
\end{tcolorbox}
\section[\texorpdfstring{\zh{特} (tè)}{特 (tè)}]{\zh{特} — the cow piece, then two pieces stacked}

\label{lesson:ZH-W450-te-special}



\begin{quote}
The word first, then the character inside it.

\begin{itemize}
\raggedright
  \item \emph{Recall:} say \emph{tèdiǎn} — \textbf{R1}, one lesson back, before it decays
\end{itemize}

\begin{quote}
\zh{对} \zh{地}
\end{quote}

\emph{Read it:} each one aloud, with nothing said first — \emph{duì}, \emph{dì}
\end{quote}

\subsection*{Script — the cow piece, then two pieces stacked}
\hlblockfigure{\detokenize{figures/ZH-W450-te-special-filmstrip.pdf}}{How \zh{特} is written, stroke by stroke}

\begin{quote}
\zh{特}
\end{quote}

\textbf{Two pieces, left before right.} Count the pieces, not the strokes.

\textbf{The left piece, the cow piece, four strokes.} A short stroke falling to the left. A bar. A long vertical down through the bar. Then a stroke that rises to the right.

\textbf{The right piece, six strokes, in two parts.} On top, the earth piece of \zh{地} with its last bar long: a short bar, a vertical, and a long bar under them. Under it, the right piece of \zh{对}: a bar, a vertical with a hook, and a dot.

4 + 6 = ten.

\subsection*{Your turn}
\begin{itemize}
\raggedright
  \item \emph{Write it:} \zh{特}, all ten, with the model in front of you
  \item \emph{Cover:} write it again from memory, and count the pieces, not the strokes
\end{itemize}

\begin{tcolorbox}[breakable,colback=teal!4,colframe=teal!35!black,title={Before you move on}]
How many strokes in \zh{特}? (\textbf{Ten}.) What is the last stroke of the left piece? (\textbf{A stroke that rises to the right}.) What sits at the foot of the right piece? (\textbf{The right piece of \zh{对}}.)

Source: \href{https://raw.githubusercontent.com/chanind/hanzi-writer-data/68d10a4b21150cae5e1ebbd223eed289cf32d90c/data/\%E7\%89\%B9.json}{Hanzi Writer Data \zh{特}, pinned stroke paths and medians}.
\end{tcolorbox}
\section[\texorpdfstring{\zh{特意} (tèyì)}{特意 (tèyì)}]{\zh{特意} (tèyì) — specially; for that very purpose}

\label{lesson:ZH-C450-teyi}



\begin{quote}
Before the new one: say the Mandarin for seafood, then the Mandarin for a distinguishing trait.
\end{quote}

\subsection*{You'll want to know: \zh{特意}}
\textbf{\zh{特意}} — \emph{tèyì} — \textquotedblleft{}specially; for that very purpose\textquotedblright{}.

\begin{quote}
\textbf{\zh{我们特意去吃这儿的风味小吃}} \emph{wǒmen tèyì qù chī zhèr de fēngwèi xiǎochī} — We went specially to try the local snacks here.
\end{quote}

\begin{cousinweb}
\textbf{\zh{特意}} (\emph{tèyì}) is built from \zh{特} 'single' and \zh{意} 'an intention', both inherited from Middle Chinese, \emph{dok} and \emph{'iH} (Guangyun, in Baxter's transcription), which glosses the second 'the will'.
\end{cousinweb}

\begin{grammarlens}[title={What you've built}]
One more word for getting around.
\end{grammarlens}

\subsection*{Your turn}
\begin{itemize}
\raggedright
  \item \emph{Say it:} \emph{tèyì}
  \item \emph{Say it:} \emph{tèyì}, once more
  \item \emph{Say it:} \emph{wǒmen tèyì qù chī zhèr de fēngwèi xiǎochī} — We went specially to try the local snacks here.
  \item \emph{Recall:} say the Mandarin for seafood, then the Mandarin for a distinguishing trait, then say \emph{tèyì} again
\end{itemize}

\begin{tcolorbox}[breakable,colback=teal!4,colframe=teal!35!black,title={Before you move on}]
What does \zh{特意} mean? (\textquotedblleft{}Specially; for that very purpose\textquotedblright{}.) Say it once more.
\end{tcolorbox}
\section[zhèngzōng]{zhèngzōng — authentic; the real thing (of food)}

\label{lesson:ZH-C450-zhengzong}



\begin{quote}
Before the new one: say the Mandarin for a distinguishing trait, then the Mandarin for specially for that reason.
\end{quote}

\subsection*{You'll want to know: zhèngzōng}
\emph{zhèngzōng} — \textquotedblleft{}authentic; the real thing (of food)\textquotedblright{}.

\begin{quote}
\emph{zhè jiā fànguǎn de cài hěn zhèngzōng} — The food at this restaurant is authentic.
\end{quote}

Its characters come later; for now it is a word for the ear.

\begin{cousinweb}
\emph{zhèngzōng} is built from \zh{正} 'right' and \emph{zōng} 'the source', both inherited from Middle Chinese, \emph{tsyengH} and \emph{tsowng} (Guangyun, in Baxter's transcription), which glosses the second 'the origin' and 'to honour'.
\end{cousinweb}

\begin{grammarlens}[title={What you've built}]
One more word for getting around.
\end{grammarlens}

\subsection*{Your turn}
\begin{itemize}
\raggedright
  \item \emph{Say it:} \emph{zhèngzōng}
  \item \emph{Say it:} \emph{zhèngzōng}, once more
  \item \emph{Say it:} \emph{zhè jiā fànguǎn de cài hěn zhèngzōng} — The food at this restaurant is authentic.
  \item \emph{Recall:} say the Mandarin for a distinguishing trait, then the Mandarin for specially for that reason, then say \emph{zhèngzōng} again
\end{itemize}

\begin{tcolorbox}[breakable,colback=teal!4,colframe=teal!35!black,title={Before you move on}]
What does zhèngzōng mean? (\textquotedblleft{}Authentic; the real thing (of food)\textquotedblright{}.) Say it once more.
\end{tcolorbox}
\section[zhāopáicài]{zhāopáicài — a restaurant's signature dish}

\label{lesson:ZH-C450-zhaopaicai}



\begin{quote}
Before the new one: say the Mandarin for specially for that reason, then the Mandarin for authentic.
\end{quote}

\subsection*{You'll want to know: zhāopáicài}
\emph{zhāopáicài} — \textquotedblleft{}a restaurant's signature dish\textquotedblright{}.

\begin{quote}
\emph{zhè shì zhè jiā fànguǎn de zhāopáicài} — This is the restaurant's signature dish.
\end{quote}

Its characters come later; for now it is a word for the ear.

\begin{cousinweb}
\emph{zhāopáicài} is built from \emph{zhāopái} 'a shop sign' and \zh{菜} 'a dish', all inherited from Middle Chinese, \emph{tsyew}, \emph{bea} and \emph{tshojH} (Guangyun, in Baxter's transcription), which glosses the first 'to beckon' and the second as a board with writing on it.
\end{cousinweb}

\begin{grammarlens}[title={What you've built}]
One more word for getting around.
\end{grammarlens}

\subsection*{Your turn}
\begin{itemize}
\raggedright
  \item \emph{Say it:} \emph{zhāopáicài}
  \item \emph{Say it:} \emph{zhāopáicài}, once more
  \item \emph{Say it:} \emph{zhè shì zhè jiā fànguǎn de zhāopáicài} — This is the restaurant's signature dish.
  \item \emph{Recall:} say the Mandarin for specially for that reason, then the Mandarin for authentic, then say \emph{zhāopáicài} again
\end{itemize}

\begin{tcolorbox}[breakable,colback=teal!4,colframe=teal!35!black,title={Before you move on}]
What does zhāopáicài mean? (\textquotedblleft{}A restaurant's signature dish\textquotedblright{}.) Say it once more.
\end{tcolorbox}
